\documentclass[12pt]{article}

%------------- MH5200-STYLE PREAMBLE (compatible subset) -------------
\usepackage[margin=1in]{geometry}
\usepackage{amsmath,amssymb,mathtools}
\usepackage{enumitem}
\usepackage{bm}
\usepackage{hyperref}
\usepackage{fancyhdr}
\usepackage{draftwatermark}
\usepackage{xcolor}

%------------- TITLE AND METADATA (REQUIRED STRUCTURE) -------------
\newcommand{\CourseCode}{MH1201}
\newcommand{\CourseName}{Linear Algebra II}
\newcommand{\DocType}{Tutorial 10 (Week 12 \& 13) -- Problems \& Solutions, Advanced Methods Edition}


\title{
\CourseCode\ \CourseName \\[0.3em]
\large \DocType
}
\author{
  Academic Year 2025/2026, Semester 2 \\[0.25em]
  \textit{Quantitative Research Society @NTU}
}
\date{\today}

%------------- HEADER & FOOTER (for page 2 onward) -------------
\fancypagestyle{main}{
  \fancyhf{}
  \fancyhead[L]{\CourseCode\ \CourseName}
  \fancyhead[R]{\href{https://qrsntu.org}{QRS@NTU}}
  \fancyfoot[C]{\thepage}
  \renewcommand{\headrulewidth}{0.4pt}
  \renewcommand{\footrulewidth}{0pt}
}

%------------- SIMPLE FOOTER (page 1 only) -------------
\fancypagestyle{plainfirst}{
  \fancyhf{}
  \renewcommand{\headrulewidth}{0pt}
  \renewcommand{\footrulewidth}{0pt}
  \fancyfoot[C]{\scriptsize\textcolor{gray}{\textit{For academic reference only. This document is provided for educational purposes and does not constitute an official question paper, answer key, or any formally authorised academic material. Its content is not guaranteed to be complete or accurate.}}}
}

%------------- WATERMARK (MH5200-style approach) -------------
\SetWatermarkText{Unofficial -- QRS@NTU -- qrsntu.org}
\SetWatermarkScale{1}
\SetWatermarkLightness{0.99}

%------------- COMMON MACROS -------------
\newcommand{\df}{\coloneqq}
\newcommand{\R}{\mathbb{R}}
\newcommand{\C}{\mathbb{C}}
\newcommand{\cL}{\mathcal{L}}
\newcommand{\cP}{\mathcal{P}}
\DeclareMathOperator{\Span}{span}
\DeclareMathOperator{\Rank}{rank}
\DeclareMathOperator{\Nullity}{nullity}
\DeclareMathOperator{\Range}{range}
\DeclareMathOperator{\Null}{Null}
\DeclareMathOperator{\Row}{Row}
\DeclareMathOperator{\tr}{trace}
\DeclareMathOperator{\proj}{proj}
\DeclareMathOperator{\diag}{diag}
\newcommand{\Id}{\mathrm{id}}
\newcommand{\ip}[2]{\left\langle #1,#2\right\rangle}
%----------------------------------------
% Optional: tighter list defaults (feel free to adjust)
\setlist[enumerate]{itemsep=0.3em, topsep=0.3em}
\setlist[itemize]{itemsep=0.2em, topsep=0.2em}

\setcounter{secnumdepth}{-1} % Globally removes numbers from all sections
\setcounter{tocdepth}{3}     % Ensures \subsubsection level shows in TOC

\begin{document}

%------------- COVER PAGE -------------
\maketitle
\thispagestyle{plainfirst}
\pagestyle{main}

%====================================================
% OVERVIEW (PAGE 1)
%====================================================
\begin{center}
  \rule{0.8\textwidth}{0.4pt}\\[4pt]
  \textbf{Overview}\\[2pt]
  \rule{0.8\textwidth}{0.4pt}
\end{center}

\noindent
This tutorial emphasizes:
\begin{itemize}[leftmargin=*]
  \item Gram--Schmidt in abstract inner product spaces and stability on orthogonal inputs;
  \item producing orthonormal bases that preserve a $T$-invariant flag, hence preserving upper-triangularity;
  \item orthogonal polynomials on $[-1,1]$ and Fourier coefficients in an orthonormal basis;
  \item inner products defined by Hermitian positive definite matrices and complex Gram--Schmidt;
  \item orthogonal complements: $W\subseteq W^{\perp\perp}$ always, and $W=W^{\perp\perp}$ in finite dimensions;
  \item explicit orthogonal decomposition $x=a+b$ with $a\in V$, $b\in V^\perp$ in $\R^4$;
  \item Frobenius inner product on $\R^{n,n}$ and describing $D^\perp$ for diagonal $D$.
\end{itemize}

\newpage
\tableofcontents
\newpage

%====================================================
% QUESTION 1
%====================================================
\section{Problem 1 (Gram--Schmidt on an orthogonal basis)}
\subsection{Problem}
Let $(V,\ip{\cdot}{\cdot})$ be a finite-dimensional inner product space, with $\dim(V)=n\ge 1$.
Let $\beta=(b_1,\dots,b_n)$ be an orthogonal ordered basis for $(V,\ip{\cdot}{\cdot})$.
Show that applying Gram--Schmidt to $\beta$ with respect to $(V,\ip{\cdot}{\cdot})$ yields $\beta$ back.

\subsection{Solution}

\subsubsection{Method 1: Projections vanish because of orthogonality}
Recall the (non-normalized) Gram--Schmidt construction:
\[
u_1=b_1,\qquad
u_k \df b_k-\sum_{j=1}^{k-1}\proj_{u_j}(b_k),
\quad
\proj_{u_j}(b_k)\df \frac{\ip{b_k}{u_j}}{\ip{u_j}{u_j}}u_j,\qquad 2\le k\le n.
\]
We prove by induction that $u_k=b_k$ for all $k$.

For $k=1$, $u_1=b_1$ is immediate.

Assume $u_j=b_j$ for all $j<k$. Since $\beta$ is orthogonal, for each $j<k$,
\[
\ip{b_k}{u_j}=\ip{b_k}{b_j}=0.
\]
Hence each projection term is $0$, so
\[
u_k=b_k-\sum_{j=1}^{k-1} 0=b_k.
\]
Thus Gram--Schmidt returns the same ordered orthogonal basis $\beta$.

\[
\boxed{\,\text{If }\beta\text{ is already orthogonal, Gram--Schmidt outputs }\beta\text{ unchanged.}\,}
\]

\subsubsection{Method 2: Uniqueness from nested spans}
In Gram--Schmidt, the output vectors satisfy:
\[
\begin{aligned}
u_k&\in \Span(b_1,\dots,b_k),\\
\Span(u_1,\dots,u_k)&=\Span(b_1,\dots,b_k),\\
\ip{u_k}{u_j}&=0\quad (j<k).
\end{aligned}
\]
Because $\beta$ is orthogonal, the vectors in $\Span(b_1,\dots,b_k)$ that are
orthogonal to $\Span(b_1,\dots,b_{k-1})$ are precisely the scalar multiples of
$b_k$. Gram--Schmidt subtracts only a vector from the preceding span, so the
$b_k$-component remains $1$. Hence $u_k=b_k$.



\subsubsection{Method 3: Orthogonal projection onto the previous flag}
For each $k$, define the previous flag subspace
\[
F_{k-1}\df \Span(b_1,\dots,b_{k-1}),
\qquad F_0\df \{0\}.
\]
The Gram--Schmidt step may be written as
\[
u_k=b_k-\proj_{F_{k-1}}(b_k),
\]
where $\proj_{F_{k-1}}(b_k)$ denotes the orthogonal projection of $b_k$ onto $F_{k-1}$.
Since $(b_1,\dots,b_n)$ is orthogonal, for each $j<k$,
\[
\ip{b_j}{b_k}=0.
\]
Therefore $b_k$ is orthogonal to every vector in $F_{k-1}$, so
\[
\proj_{F_{k-1}}(b_k)=0.
\]
Hence
\[
u_k=b_k.
\]
Thus every Gram--Schmidt step leaves the input vector unchanged.

\[
\boxed{\,u_k=b_k\text{ for every }k=1,\dots,n.\,}
\]

\paragraph{Convention note.}
This conclusion is for the non-normalized Gram--Schmidt algorithm. If one's definition of Gram--Schmidt includes normalization at every step, then the output is instead
\[
\left(\frac{b_1}{\|b_1\|},\dots,\frac{b_n}{\|b_n\|}\right).
\]
Thus the precise statement depends on whether ``Gram--Schmidt'' means orthogonalisation or orthonormalisation.


%====================================================
% QUESTION 2
%====================================================
\newpage
\section{Problem 2 (Preserving upper-triangularity)}
\subsection{Problem}
Let
\[
\beta \df ((1,0,0),(1,1,1),(1,1,2)),
\]
which is an ordered basis for $\R^3$ (with the standard inner product $\cdot$).
Let $T\in \cL(\R^3)$, and suppose that $[T]_{\beta}^{\beta}\in \R^{3,3}$ is upper triangular.
Derive an orthonormal ordered basis $\gamma$ for $(\R^3,\cdot)$ for which $[T]_{\gamma}^{\gamma}\in\R^{3,3}$ is upper triangular.

\subsection{Solution}

\subsubsection{Method 1: Gram--Schmidt preserves the flag; hence upper triangular}
Write $\beta=(b_1,b_2,b_3)$ with
\[
b_1=(1,0,0),\quad b_2=(1,1,1),\quad b_3=(1,1,2).
\]
Since $[T]_{\beta}^{\beta}$ is upper triangular, we have the equivalent $T$-invariance of the coordinate flag:
\[
\begin{aligned}
T\bigl(\Span(b_1)\bigr)&\subseteq \Span(b_1),\\
T\bigl(\Span(b_1,b_2)\bigr)&\subseteq \Span(b_1,b_2),\\
T\bigl(\Span(b_1,b_2,b_3)\bigr)&\subseteq \Span(b_1,b_2,b_3)=\R^3.
\end{aligned}
\]
Apply Gram--Schmidt to $\beta$ in the standard inner product.

\medskip
\noindent\textbf{Step 1: Compute $\gamma$ by Gram--Schmidt.}
Let
\[
u_1=b_1,\qquad
u_2=b_2-\proj_{u_1}(b_2),\qquad
u_3=b_3-\proj_{u_1}(b_3)-\proj_{u_2}(b_3).
\]
Compute:
\[
\proj_{u_1}(b_2)=\frac{b_2\cdot u_1}{u_1\cdot u_1}u_1=\frac{1}{1}b_1=b_1,
\quad\Rightarrow\quad u_2=b_2-b_1=(0,1,1).
\]
Similarly,
\[
\proj_{u_1}(b_3)=\frac{b_3\cdot u_1}{u_1\cdot u_1}u_1=b_1,
\quad\Rightarrow\quad b_3-\proj_{u_1}(b_3)=(0,1,2).
\]
Now
\[
\proj_{u_2}(b_3)=\proj_{u_2}\bigl((0,1,2)\bigr)
=\frac{(0,1,2)\cdot (0,1,1)}{(0,1,1)\cdot(0,1,1)}(0,1,1)
=\frac{3}{2}(0,1,1).
\]
Thus
\[
u_3=(0,1,2)-\frac{3}{2}(0,1,1)=\left(0,-\frac12,\frac12\right).
\]
Normalize:
\[
\|u_1\|=1,\quad \|u_2\|=\sqrt2,\quad \|u_3\|=\sqrt{\frac12}.
\]
Hence an orthonormal basis is
\[
\gamma=(g_1,g_2,g_3)\df
\left(
(1,0,0),\ \left(0,\frac{1}{\sqrt2},\frac{1}{\sqrt2}\right),\ \left(0,-\frac{1}{\sqrt2},\frac{1}{\sqrt2}\right)
\right).
\]
\[
\boxed{\,
\gamma=\left(
(1,0,0),\ \left(0,\frac{1}{\sqrt2},\frac{1}{\sqrt2}\right),\ \left(0,-\frac{1}{\sqrt2},\frac{1}{\sqrt2}\right)
\right)\ \text{is orthonormal.}
\,}
\]

\medskip
\noindent\textbf{Step 2: Show $[T]_\gamma^\gamma$ is upper triangular.}
A key Gram--Schmidt property is
\[
g_k\in \Span(b_1,\dots,b_k)\qquad (k=1,2,3).
\]
Therefore
\[
\Span(g_1,\dots,g_k)=\Span(b_1,\dots,b_k)\qquad (k=1,2,3).
\]
Since $T(\Span(b_1,\dots,b_k))\subseteq \Span(b_1,\dots,b_k)$, we get
\[
T(\Span(g_1,\dots,g_k))\subseteq \Span(g_1,\dots,g_k)\qquad (k=1,2,3).
\]
This is exactly the condition that the matrix of $T$ in the ordered basis $\gamma$ is upper triangular.

\[
\boxed{\, [T]_{\gamma}^{\gamma}\ \text{is upper triangular.}\,}
\]

\subsubsection{Method 2: Upper-triangular similarity via an upper-triangular change-of-basis matrix}
Let $B=[T]_{\beta}^{\beta}$ (upper triangular). Let $P$ be the change-of-basis matrix from $\gamma$-coordinates to $\beta$-coordinates:
\[
P \df [\Id]_{\beta}^{\gamma}.
\]
Because Gram--Schmidt produces $g_k\in\Span(b_1,\dots,b_k)$, the coordinate vector $[g_k]_\beta$ has no components in $b_{k+1},\dots,b_3$,
so $P$ is upper triangular with nonzero diagonal (hence invertible).
Then
\[
[T]_{\gamma}^{\gamma}=P^{-1}\,[T]_{\beta}^{\beta}\,P=P^{-1}BP.
\]
A product (and inverse) of upper-triangular invertible matrices is upper triangular, hence $[T]_{\gamma}^{\gamma}$ is upper triangular.



\subsubsection{Method 3: Flag-slice construction without full Gram--Schmidt}
Define the flag associated with the original ordered basis:
\[
F_k\df \Span(b_1,\dots,b_k),\qquad k=1,2,3.
\]
Because $[T]_{\beta}^{\beta}$ is upper triangular, the flag is $T$-invariant:
\[
T(F_k)\subseteq F_k\qquad (k=1,2,3).
\]
Instead of carrying out Gram--Schmidt mechanically, construct the orthonormal basis by taking the new unit direction in each orthogonal slice:
\[
g_k\in F_k\cap F_{k-1}^{\perp},
\qquad \|g_k\|=1,
\qquad F_0\df\{0\}.
\]
For this problem,
\[
F_1=\Span\{(1,0,0)\},
\]
so
\[
g_1=(1,0,0).
\]
Next,
\[
F_2=\Span\{(1,0,0),(1,1,1)\}=\{(a,b,b):a,b\in\R\}.
\]
The vector in $F_2$ perpendicular to $g_1$ must have first coordinate $0$, hence is proportional to $(0,1,1)$. Therefore
\[
g_2=\frac{1}{\sqrt2}(0,1,1).
\]
Finally, since $F_3=\R^3$, the last direction is the unit vector perpendicular to $F_2$. A vector perpendicular to both $(1,0,0)$ and $(0,1,1)$ is proportional to $(0,-1,1)$, so we may take
\[
g_3=\frac{1}{\sqrt2}(0,-1,1).
\]
Thus
\[
\boxed{\,
\gamma=\left((1,0,0),\frac{1}{\sqrt2}(0,1,1),\frac{1}{\sqrt2}(0,-1,1)\right).
\,}
\]
Moreover,
\[
\Span(g_1,\dots,g_k)=F_k\qquad(k=1,2,3).
\]
Since each $F_k$ is $T$-invariant, the ordered basis $\gamma$ also gives an invariant flag. Hence
\[
\boxed{\,[T]_{\gamma}^{\gamma}\text{ is upper triangular.}\,}
\]
This method is usually faster than computing all projection coefficients.

\subsubsection{Method 4: QR-factorisation and triangular similarity}
Let
\[
C\df [b_1\ b_2\ b_3]
=\begin{bmatrix}
1&1&1\\
0&1&1\\
0&1&2
\end{bmatrix}.
\]
Applying Gram--Schmidt to the columns of $C$ gives a QR factorisation
\[
C=QR,
\]
where the columns of $Q$ form the orthonormal basis $\gamma$ and $R$ is invertible upper triangular.

Let
\[
B\df [T]_{\beta}^{\beta}.
\]
By hypothesis, $B$ is upper triangular. Since $R$ is an upper-triangular change-of-coordinates matrix between the two flag-adapted bases, the matrix of $T$ in the orthonormal basis is
\[
[T]_{\gamma}^{\gamma}=RBR^{-1}.
\]
The inverse of an invertible upper-triangular matrix is upper triangular, and products of upper-triangular matrices are upper triangular. Therefore
\[
\boxed{\,[T]_{\gamma}^{\gamma}\text{ is upper triangular.}\,}
\]
This method packages the whole argument into the fact that Gram--Schmidt is exactly QR factorisation.


%====================================================
% QUESTION 3
%====================================================
\newpage
\section{Problem 3 (Orthogonal polynomials on $[-1,1]$)}
\subsection{Problem}
Define a function $\ip{\cdot}{\cdot}:\cP_3(\R)\times \cP_3(\R)\to \R$ as follows:
\[
\forall P,Q\in\cP_3(\R):\qquad
\ip{P}{Q}\df \int_{-1}^{1} P(x)\,Q(x)\,dx.
\]
\begin{enumerate}[label=(\alph*)]
\item Show that $\ip{\cdot}{\cdot}$ is an inner product on $\cP_3(\R)$.
\item Apply Gram--Schmidt to $(1,X,X^2,X^3)$ with respect to $(\cP_3(\R),\ip{\cdot}{\cdot})$.
\item Normalize your answer to Part (b) to obtain an orthonormal basis $\beta$ for $(\cP_3(\R),\ip{\cdot}{\cdot})$.
\item Compute the $\beta$-Fourier coefficients of $1-X+X^2-X^3$ with respect to $(\cP_3(\R),\ip{\cdot}{\cdot})$.
\end{enumerate}

\subsection{Solution}

\subsubsection{Method 1: Direct axiom verification + explicit Gram--Schmidt computations}
\begin{enumerate}[label=(\alph*)]
\item \textbf{Inner product axioms.}
Bilinearity follows from linearity of the integral:
\[
\ip{aP_1+bP_2}{Q}=a\ip{P_1}{Q}+b\ip{P_2}{Q},\qquad
\ip{P}{aQ_1+bQ_2}=a\ip{P}{Q_1}+b\ip{P}{Q_2}.
\]
Symmetry holds because $P(x)Q(x)=Q(x)P(x)$:
\[
\ip{P}{Q}=\int_{-1}^1 P Q=\int_{-1}^1 Q P=\ip{Q}{P}.
\]
Positive definiteness: for $P\neq 0$,
\[
\ip{P}{P}=\int_{-1}^1 P(x)^2\,dx>0.
\]
Indeed, $P(x)^2\ge 0$ and is continuous; if the integral were $0$, then $P(x)^2=0$ for all $x\in[-1,1]$, so $P$ vanishes on an interval,
hence $P$ is the zero polynomial, contradiction. Therefore $\ip{\cdot}{\cdot}$ is an inner product.

\[
\boxed{\,\ip{\cdot}{\cdot}\text{ is an inner product on }\cP_3(\R).\,}
\]

\item \textbf{Gram--Schmidt on $(1,X,X^2,X^3)$.}
Let $v_0=1$, $v_1=X$, $v_2=X^2$, $v_3=X^3$.
Define
\[
\begin{aligned}
u_0&=v_0, & u_1&=v_1-\proj_{u_0}(v_1),\\
u_2&=v_2-\proj_{u_0}(v_2)-\proj_{u_1}(v_2), &
u_3&=v_3-\sum_{j=0}^{2}\proj_{u_j}(v_3).
\end{aligned}
\]

Compute required inner products using parity:
\[
\int_{-1}^1 x\,dx=0,\quad \int_{-1}^1 x^3\,dx=0,\quad \int_{-1}^1 x^5\,dx=0.
\]
Also
\[
\int_{-1}^1 1\,dx=2,\quad \int_{-1}^1 x^2\,dx=\frac{2}{3},\quad \int_{-1}^1 x^4\,dx=\frac{2}{5},\quad \int_{-1}^1 x^6\,dx=\frac{2}{7}.
\]

\underline{Step 1: $u_0=1$.}

\underline{Step 2: $u_1$.}
Since $\ip{X}{1}=\int_{-1}^1 x\,dx=0$, we get
\[
u_1=X.
\]

\underline{Step 3: $u_2$.}
Compute
\[
\proj_{u_0}(X^2)=\frac{\ip{X^2}{1}}{\ip{1}{1}}1=\frac{\frac{2}{3}}{2}\cdot 1=\frac{1}{3}.
\]
Also $\ip{X^2}{X}=\int_{-1}^1 x^3\,dx=0$, so $\proj_{u_1}(X^2)=0$.
Hence
\[
u_2=X^2-\frac{1}{3}.
\]

\underline{Step 4: $u_3$.}
We have $\ip{X^3}{1}=\int_{-1}^1 x^3\,dx=0$, and
\[
\proj_{u_1}(X^3)=\frac{\ip{X^3}{X}}{\ip{X}{X}}X
=\frac{\int_{-1}^1 x^4\,dx}{\int_{-1}^1 x^2\,dx}X
=\frac{\frac{2}{5}}{\frac{2}{3}}X=\frac{3}{5}X.
\]
Finally,
\[
\ip{X^3}{u_2}=\int_{-1}^1 x^3\left(x^2-\frac13\right)\,dx
=\int_{-1}^1 \left(x^5-\frac13 x^3\right)\,dx=0,
\]
so $\proj_{u_2}(X^3)=0$. Therefore
\[
u_3=X^3-\frac{3}{5}X.
\]

Thus Gram--Schmidt yields the orthogonal list
\[
\boxed{\,\bigl(1,\ X,\ X^2-\tfrac13,\ X^3-\tfrac35 X\bigr).\,}
\]

\item \textbf{Normalize to an orthonormal basis $\beta$.}
Compute norms:

\underline{$u_0=1$:} $\|u_0\|^2=\int_{-1}^1 1\,dx=2$, so $\|u_0\|=\sqrt2$.

\underline{$u_1=X$:} $\|u_1\|^2=\int_{-1}^1 x^2\,dx=\frac{2}{3}$, so $\|u_1\|=\sqrt{\frac{2}{3}}$.

\underline{$u_2=X^2-\frac13$:}
\[
\|u_2\|^2=\int_{-1}^1 \left(x^2-\frac13\right)^2dx
=\int_{-1}^1\left(x^4-\frac{2}{3}x^2+\frac{1}{9}\right)dx
=\frac{2}{5}-\frac{2}{3}\cdot\frac{2}{3}+\frac{1}{9}\cdot 2
=\frac{8}{45}.
\]
So $\|u_2\|=\sqrt{\frac{8}{45}}$.

\underline{$u_3=X^3-\frac35 X$:}
\[
\|u_3\|^2=\int_{-1}^1\left(x^3-\frac{3}{5}x\right)^2dx
=\int_{-1}^1\left(x^6-\frac{6}{5}x^4+\frac{9}{25}x^2\right)dx
=\frac{2}{7}-\frac{6}{5}\cdot\frac{2}{5}+\frac{9}{25}\cdot\frac{2}{3}
=\frac{8}{175}.
\]
So $\|u_3\|=\sqrt{\frac{8}{175}}$.

Define the orthonormal basis $\beta=(\beta_0,\beta_1,\beta_2,\beta_3)$ by $\beta_j=u_j/\|u_j\|$:
\[
\beta_0=\frac{1}{\sqrt2},\qquad
\beta_1=\sqrt{\frac{3}{2}}\,X,\qquad
\beta_2=\sqrt{\frac{45}{8}}\left(X^2-\frac13\right),\qquad
\beta_3=\sqrt{\frac{175}{8}}\left(X^3-\frac35 X\right).
\]
\[
\boxed{\,\beta=\left(\frac{1}{\sqrt2},\ \sqrt{\frac{3}{2}}X,\ \sqrt{\frac{45}{8}}\left(X^2-\frac13\right),\ \sqrt{\frac{175}{8}}\left(X^3-\frac35 X\right)\right)\ \text{is orthonormal.}\,}
\]

\item \textbf{$\beta$-Fourier coefficients of $f(X)=1-X+X^2-X^3$.}
Since $\beta$ is orthonormal, the Fourier coefficient of $f$ along $\beta_j$ is
\[
c_j=\ip{f}{\beta_j}.
\]

\underline{$c_0$:}
\[
c_0=\ip{f}{\tfrac{1}{\sqrt2}}=\frac{1}{\sqrt2}\int_{-1}^1 (1-x+x^2-x^3)\,dx
=\frac{1}{\sqrt2}\left(2+\frac{2}{3}\right)=\frac{8}{3\sqrt2}.
\]

\underline{$c_1$:}
\[
\begin{aligned}
c_1&=\ip{f}{\sqrt{\frac32}x}
=\sqrt{\frac32}\int_{-1}^1 (1-x+x^2-x^3)x\,dx\\
&=\sqrt{\frac32}\int_{-1}^1(x-x^2+x^3-x^4)\,dx\\
&=\sqrt{\frac32}\left(-\frac{2}{3}-\frac{2}{5}\right)
=-\frac{16}{15}\sqrt{\frac32}.
\end{aligned}
\]

\underline{$c_2$:}
\[
c_2=\ip{f}{\sqrt{\frac{45}{8}}\left(x^2-\frac13\right)}
=\sqrt{\frac{45}{8}}\int_{-1}^1 (1-x+x^2-x^3)\left(x^2-\frac13\right)\,dx.
\]
Expanding and using parity (odd terms integrate to $0$) gives
\[
\int_{-1}^1 (1-x+x^2-x^3)\left(x^2-\frac13\right)\,dx
=\frac{8}{45}.
\]
Thus
\[
c_2=\sqrt{\frac{45}{8}}\cdot \frac{8}{45}=\frac{4\sqrt5}{15\sqrt2}.
\]

\underline{$c_3$:}
\[
c_3=\ip{f}{\sqrt{\frac{175}{8}}\left(x^3-\frac35 x\right)}
=\sqrt{\frac{175}{8}}\int_{-1}^1 (1-x+x^2-x^3)\left(x^3-\frac35 x\right)\,dx.
\]
Keeping only even powers after expansion yields
\[
\int_{-1}^1 (1-x+x^2-x^3)\left(x^3-\frac35 x\right)\,dx
=-\frac{8}{175}.
\]
Hence
\[
c_3=\sqrt{\frac{175}{8}}\cdot\left(-\frac{8}{175}\right)=-\frac{4\sqrt7}{35\sqrt2}.
\]

Therefore the $\beta$-Fourier coefficients are
\[
\boxed{\,
(c_0,c_1,c_2,c_3)=\left(\frac{8}{3\sqrt2},\ -\frac{16}{15}\sqrt{\frac32},\ \frac{4\sqrt5}{15\sqrt2},\ -\frac{4\sqrt7}{35\sqrt2}\right).
\,}
\]
\end{enumerate}


\newpage
\subsubsection{Method 2: Legendre-polynomial shortcut}
\begin{enumerate}[label=(\alph*)]
\item \textbf{Inner product axioms.}


For \(a,b\in\R\) and \(P_1,P_2,Q\in\cP_3(\R)\),
\[
\ip{aP_1+bP_2}{Q}
=\int_{-1}^{1}(aP_1(x)+bP_2(x))Q(x)\,dx
=a\ip{P_1}{Q}+b\ip{P_2}{Q}.
\]
Similarly, by linearity in the second argument,
\[
\ip{P}{aQ_1+bQ_2}
=a\ip{P}{Q_1}+b\ip{P}{Q_2}.
\]
Also,
\[
\ip{P}{Q}
=\int_{-1}^{1}P(x)Q(x)\,dx
=\int_{-1}^{1}Q(x)P(x)\,dx
=\ip{Q}{P}.
\]
Finally,
\[
\ip{P}{P}
=\int_{-1}^{1}P(x)^2\,dx\ge 0.
\]
If \(\ip{P}{P}=0\), then
\[
\int_{-1}^{1}P(x)^2\,dx=0.
\]
Since \(P(x)^2\ge 0\) is continuous on \([-1,1]\), this implies
\[
P(x)^2=0
\qquad\forall x\in[-1,1].
\]
Hence \(P(x)=0\) for all \(x\in[-1,1]\). Since a nonzero polynomial cannot vanish on an interval, we must have \(P=0\). Therefore
\[
\ip{P}{P}=0 \iff P=0.
\]
So \(\ip{\cdot}{\cdot}\) is an inner product on \(\cP_3(\R)\).

\[
\boxed{\ip{\cdot}{\cdot}\text{ is an inner product on }\cP_3(\R).}
\]

\item \textbf{Apply Gram--Schmidt using Legendre polynomials.}

The standard Legendre polynomials are
\[
L_0(x)=1,\qquad
L_1(x)=x,\qquad
L_2(x)=\frac12(3x^2-1),\qquad
L_3(x)=\frac12(5x^3-3x).
\]
They satisfy the orthogonality relation
\[
\int_{-1}^{1}L_m(x)L_n(x)\,dx=0
\qquad(m\neq n).
\]
Therefore
\[
(L_0,L_1,L_2,L_3)
\]
is an orthogonal list in \(\cP_3(\R)\).

We now connect this with Gram--Schmidt applied to
\[
(1,X,X^2,X^3).
\]
Gram--Schmidt produces an orthogonal polynomial of degree \(0\), then degree \(1\), then degree \(2\), then degree \(3\). Since \(L_n\) is also an orthogonal polynomial of degree \(n\), the Gram--Schmidt output must agree with the Legendre polynomial \(L_n\) up to a nonzero scalar multiple.

Now rewrite \(L_2\) and \(L_3\) in monic-like form:
\[
L_2(X)=\frac12(3X^2-1)
=\frac32\left(X^2-\frac13\right),
\]
so
\[
X^2-\frac13=\frac23L_2(X).
\]
Similarly,
\[
L_3(X)=\frac12(5X^3-3X)
=\frac52\left(X^3-\frac35X\right),
\]
so
\[
X^3-\frac35X=\frac25L_3(X).
\]

Hence the Gram--Schmidt orthogonal list obtained from
\[
(1,X,X^2,X^3)
\]
is
\[
\boxed{
\left(
1,\,
X,\,
X^2-\frac13,\,
X^3-\frac35X
\right).
}
\]

\item \textbf{Normalize the orthogonal list.}

The Legendre norm formula is
\[
\int_{-1}^{1}L_n(x)^2\,dx=\frac{2}{2n+1}.
\]
Hence the normalized Legendre polynomial is
\[
\phi_n(X)=\sqrt{\frac{2n+1}{2}}L_n(X).
\]

For \(n=0,1,2,3\), this gives
\[
\phi_0(X)=\frac{1}{\sqrt2},
\]
\[
\phi_1(X)=\sqrt{\frac32}X,
\]
\[
\phi_2(X)
=\sqrt{\frac52}\cdot \frac12(3X^2-1),
\]
\[
\phi_3(X)
=\sqrt{\frac72}\cdot \frac12(5X^3-3X).
\]

Equivalently, using
\[
L_2(X)=\frac32\left(X^2-\frac13\right),
\qquad
L_3(X)=\frac52\left(X^3-\frac35X\right),
\]
we may write
\[
\phi_2(X)
=\sqrt{\frac52}\cdot \frac32\left(X^2-\frac13\right)
=\sqrt{\frac{45}{8}}\left(X^2-\frac13\right),
\]
and
\[
\phi_3(X)
=\sqrt{\frac72}\cdot \frac52\left(X^3-\frac35X\right)
=\sqrt{\frac{175}{8}}\left(X^3-\frac35X\right).
\]

Therefore an orthonormal basis is
\[
\boxed{
\beta=
\left(
\frac{1}{\sqrt2},\,
\sqrt{\frac32}X,\,
\sqrt{\frac{45}{8}}\left(X^2-\frac13\right),\,
\sqrt{\frac{175}{8}}\left(X^3-\frac35X\right)
\right).
}
\]

\item \textbf{Compute the \(\beta\)-Fourier coefficients of \(f(X)=1-X+X^2-X^3\).}

Since the basis in Part (c) is the normalized Legendre basis, it is convenient to first express \(f\) in the unnormalized Legendre basis.

We use
\[
L_0(X)=1,\qquad L_1(X)=X,
\]
and from
\[
L_2(X)=\frac12(3X^2-1),
\]
we get
\[
3X^2-1=2L_2(X),
\]
so
\[
X^2=\frac13+\frac23L_2(X).
\]
Also, from
\[
L_3(X)=\frac12(5X^3-3X),
\]
we get
\[
5X^3-3X=2L_3(X),
\]
so
\[
X^3=\frac35X+\frac25L_3(X)
=\frac35L_1(X)+\frac25L_3(X).
\]

Now substitute these into
\[
f(X)=1-X+X^2-X^3.
\]
Then
\[
\begin{aligned}
f(X)
&=L_0(X)-L_1(X)
+\left(\frac13L_0(X)+\frac23L_2(X)\right)
-\left(\frac35L_1(X)+\frac25L_3(X)\right)\\
&=\frac43L_0(X)-\frac85L_1(X)+\frac23L_2(X)-\frac25L_3(X).
\end{aligned}
\]

The normalized Legendre basis is
\[
\phi_n(X)=\sqrt{\frac{2n+1}{2}}L_n(X).
\]
Therefore
\[
L_n(X)=\sqrt{\frac{2}{2n+1}}\phi_n(X).
\]
So if
\[
f(X)=a_0L_0(X)+a_1L_1(X)+a_2L_2(X)+a_3L_3(X),
\]
then the Fourier coefficient of \(f\) along \(\phi_n\) is
\[
c_n=a_n\sqrt{\frac{2}{2n+1}}.
\]

Here,
\[
a_0=\frac43,\qquad
a_1=-\frac85,\qquad
a_2=\frac23,\qquad
a_3=-\frac25.
\]
Thus
\[
c_0=\frac43\sqrt2
=\frac{8}{3\sqrt2},
\]
\[
c_1=-\frac85\sqrt{\frac23}
=-\frac{16}{15}\sqrt{\frac32},
\]
\[
c_2=\frac23\sqrt{\frac25}
=\frac{4\sqrt5}{15\sqrt2},
\]
and
\[
c_3=-\frac25\sqrt{\frac27}
=-\frac{4\sqrt7}{35\sqrt2}.
\]

Therefore the \(\beta\)-Fourier coefficients are
\[
\boxed{
(c_0,c_1,c_2,c_3)
=
\left(
\frac{8}{3\sqrt2},\,
-\frac{16}{15}\sqrt{\frac32},\,
\frac{4\sqrt5}{15\sqrt2},\,
-\frac{4\sqrt7}{35\sqrt2}
\right).
}
\]
\end{enumerate}

\newpage
\subsubsection{Method 3: Even--odd parity block shortcut}
The inner product
\[
\ip{P}{Q}=\int_{-1}^{1}P(x)Q(x)\,dx
\]
has an important symmetry: odd functions integrate to zero over $[-1,1]$.
Therefore even polynomials are orthogonal to odd polynomials. Hence Gram--Schmidt separates into two independent blocks:
\[
\Span\{1,X^2\}\quad\text{and}\quad \Span\{X,X^3\}.
\]
For the even block,
\[
u_0=1,
\qquad
u_2=X^2-\frac{\ip{1}{X^2}}{\ip{1}{1}}1
=X^2-\frac13.
\]
For the odd block,
\[
u_1=X,
\qquad
u_3=X^3-\frac{\ip{X}{X^3}}{\ip{X}{X}}X
=X^3-\frac{3}{5}X.
\]
This recovers
\[
\boxed{\,(1,\ X,\ X^2-\tfrac13,\ X^3-\tfrac35X)\,}
\]
with much less computation. The method works whenever the interval and weight function are symmetric, so that even--odd orthogonality is available.


%====================================================
% QUESTION 4
%====================================================
\newpage
\section{Problem 4 (Complex inner product via a Hermitian positive definite matrix)}
\subsection{Problem}
Define a function $\ip{\cdot}{\cdot}:\C^{3,1}\times \C^{3,1}\to \C$ as follows:
\[
\forall u,v\in \C^{3,1}:\qquad
\ip{u}{v}\df u^{\dagger}
\begin{bmatrix}
2 & 1+i & 0\\
1-i & 3 & 2\\
0 & 2 & 4
\end{bmatrix}
v,
\]
where ${}^\dagger$ denotes conjugate transposition.
\begin{enumerate}[label=(\alph*)]
\item Show that $\ip{\cdot}{\cdot}$ is an inner product on $\C^{3,1}$.
\item Apply Gram--Schmidt to the standard ordered basis for $\C^{3,1}$ with respect to $(\C^{3,1},\ip{\cdot}{\cdot})$.
\item Normalize your answer to Part (b) to obtain an orthonormal basis $\beta$ for $(\C^{3,1},\ip{\cdot}{\cdot})$.
\item Compute the $\beta$-Fourier coefficients of $\begin{bmatrix}1\\0\\ i\end{bmatrix}$ with respect to $(\C^{3,1},\ip{\cdot}{\cdot})$.
\end{enumerate}

\subsection{Solution}
Let
\[
H\df
\begin{bmatrix}
2 & 1+i & 0\\
1-i & 3 & 2\\
0 & 2 & 4
\end{bmatrix},
\qquad
\ip{u}{v}=u^\dagger H v.
\]
We use the standard convention: conjugate-linear in the first argument and linear in the second.

\subsubsection{Method 1: Hermitian positive definite $\Rightarrow$ inner product; then compute Gram--Schmidt explicitly}
\begin{enumerate}[label=(\alph*)]
\item \textbf{Axiom 1: sesquilinearity.}
For $a,b\in\C$ and $u_1,u_2,v\in\C^{3,1}$,
\[
\ip{au_1+bu_2}{v}=(au_1+bu_2)^\dagger H v=\overline{a}\,u_1^\dagger Hv+\overline{b}\,u_2^\dagger Hv
=\overline{a}\ip{u_1}{v}+\overline{b}\ip{u_2}{v},
\]
and similarly $\ip{u}{av_1+bv_2}=a\ip{u}{v_1}+b\ip{u}{v_2}$.

\smallskip
\noindent\textbf{Axiom 2: conjugate symmetry.}
Since $H^\dagger=H$ (check immediately from entries), we have
\[
\overline{\ip{v}{u}}=\overline{v^\dagger H u}=u^\dagger H^\dagger v=u^\dagger H v=\ip{u}{v}.
\]

\smallskip
\noindent\textbf{Axiom 3: positive definiteness.}
Because $H$ is Hermitian, $\ip{u}{u}=u^\dagger H u\in\R$.
We show $u^\dagger H u>0$ for $u\neq 0$ using Sylvester's criterion (leading principal minors):
\[
\det([2])=2>0,\qquad
\det\begin{bmatrix}2&1+i\\1-i&3\end{bmatrix}
=6-(1+i)(1-i)=6-2=4>0,
\]
and
\[
\det(H)=2\det\begin{bmatrix}3&2\\2&4\end{bmatrix}-(1+i)\det\begin{bmatrix}1-i&2\\0&4\end{bmatrix}
=2(12-4)- (1+i)\cdot 4(1-i)
=16-4\cdot 2=8>0.
\]
Hence $H$ is Hermitian positive definite, so $u^\dagger H u>0$ for all $u\neq 0$.
Therefore $\ip{\cdot}{\cdot}$ is an inner product.

\[
\boxed{\,\ip{\cdot}{\cdot}\text{ is an inner product on }\C^{3,1}. \,}
\]

\item \textbf{Gram--Schmidt on the standard basis.}
Let $e_1,e_2,e_3$ be the standard basis of $\C^{3,1}$.
Define
\[
u_1=e_1,\qquad
u_2=e_2-\proj_{u_1}(e_2),\qquad
u_3=e_3-\proj_{u_1}(e_3)-\proj_{u_2}(e_3),
\]
where $\proj_{u}(v)=\dfrac{\ip{u}{v}}{\ip{u}{u}}u$.

Compute the needed inner products using $\ip{e_i}{e_j}=H_{ij}$:
\[
\ip{u_1}{u_1}=\ip{e_1}{e_1}=2,\qquad
\ip{u_1}{e_2}=\ip{e_1}{e_2}=1+i,\qquad
\ip{u_1}{e_3}=\ip{e_1}{e_3}=0.
\]
Hence
\[
u_2=e_2-\frac{1+i}{2}e_1.
\]
Next,
\[
\ip{u_2}{u_2}=2,
\qquad
\ip{u_2}{e_3}=2,
\]
so $\proj_{u_2}(e_3)=\dfrac{2}{2}u_2=u_2$ and $\proj_{u_1}(e_3)=0$.
Thus
\[
u_3=e_3-u_2=e_3-e_2+\frac{1+i}{2}e_1.
\]
So Gram--Schmidt yields the orthogonal ordered basis
\[
\boxed{\,
(u_1,u_2,u_3)=\left(e_1,\ e_2-\frac{1+i}{2}e_1,\ \frac{1+i}{2}e_1-e_2+e_3\right).
\,}
\]

\item \textbf{Normalize to an orthonormal basis $\beta$.}
We already computed $\ip{u_1}{u_1}=2$ and $\ip{u_2}{u_2}=2$.
Also,
\[
\ip{u_3}{u_3}=2
\]
(one direct check is $Hu_3=(0,0,2)^{\mathsf T}$, hence $u_3^\dagger Hu_3=2$).
Therefore $\|u_k\|=\sqrt2$ for $k=1,2,3$.

Define
\[
\beta=(\beta_1,\beta_2,\beta_3)\df \left(\frac{u_1}{\sqrt2},\ \frac{u_2}{\sqrt2},\ \frac{u_3}{\sqrt2}\right).
\]
Explicitly,
\[
\boxed{
\beta=
\left(
\frac{1}{\sqrt2}\!\begin{bmatrix}1\\0\\0\end{bmatrix},\ 
\frac{1}{\sqrt2}\!\begin{bmatrix}-\frac{1+i}{2}\\1\\0\end{bmatrix},\ 
\frac{1}{\sqrt2}\!\begin{bmatrix}\frac{1+i}{2}\\-1\\1\end{bmatrix}
\right)
\ \text{is orthonormal in }(\C^{3,1},\ip{\cdot}{\cdot}).
}
\]

\item \textbf{$\beta$-Fourier coefficients of $v=\begin{bmatrix}1\\0\\ i\end{bmatrix}$.}
Since $\beta$ is orthonormal, coefficients are
\[
c_k=\ip{\beta_k}{v}\qquad (k=1,2,3).
\]
First compute
\[
Hv=
\begin{bmatrix}
2 & 1+i & 0\\
1-i & 3 & 2\\
0 & 2 & 4
\end{bmatrix}
\begin{bmatrix}1\\0\\ i\end{bmatrix}
=
\begin{bmatrix}
2\\
1+i\\
4i
\end{bmatrix}.
\]
Now:
\[
c_1=\ip{\beta_1}{v}
=\left(\frac{1}{\sqrt2}e_1\right)^\dagger Hv
=\frac{1}{\sqrt2}\,(2)=\sqrt2.
\]
Let $\alpha\df \frac{1+i}{2}$, so $u_2=e_2-\alpha e_1$ and $\overline{\alpha}=\frac{1-i}{2}$.
Then
\[
c_2=\ip{\beta_2}{v}
=\frac{1}{\sqrt2}\ip{u_2}{v}
=\frac{1}{\sqrt2}\,u_2^\dagger(Hv)
=\frac{1}{\sqrt2}\left((-\overline{\alpha})\cdot 2 + 1\cdot (1+i)\right)
=\frac{1}{\sqrt2}(2i)=\sqrt2\,i.
\]
Similarly, $u_3=\alpha e_1-e_2+e_3$, so
\[
c_3=\ip{\beta_3}{v}
=\frac{1}{\sqrt2}\,u_3^\dagger(Hv)
=\frac{1}{\sqrt2}\left(\overline{\alpha}\cdot 2 - (1+i) + 4i\right)
=\frac{1}{\sqrt2}(2i)=\sqrt2\,i.
\]
Therefore the Fourier coefficients are
\[
\boxed{\, (c_1,c_2,c_3)=\left(\sqrt2,\ \sqrt2\,i,\ \sqrt2\,i\right).\,}
\]
\end{enumerate}

\newpage
\subsubsection{Method 2: Matrix/Gram-matrix formulation}
The theorems and definitions used in this method are the following.
\begin{itemize}[leftmargin=*]
  \item \textbf{Hermitian positive definite matrix theorem.} If $H\in\C^{n,n}$ is Hermitian positive definite, then
  \[
  \ip{u}{v}\df u^\dagger Hv
  \]
  defines an inner product on $\C^{n,1}$ under the convention that the first slot is conjugate-linear and the second slot is linear.
  \item \textbf{Gram matrix definition.} Given a basis $(e_1,\dots,e_n)$ of an inner product space, its Gram matrix is
  $G=(\ip{e_i}{e_j})_{i,j}$. For the standard basis and the inner product $\ip{u}{v}=u^\dagger Hv$, this Gram matrix is exactly $H$.
\end{itemize}

\begin{enumerate}[label=(\alph*)]
\item Since
\[
H^\dagger=H,
\]
the matrix $H$ is Hermitian. Its leading principal minors are
\[
\Delta_1=2>0,
\]
\[
\Delta_2=\det\begin{bmatrix}2&1+i\\1-i&3\end{bmatrix}
=6-(1+i)(1-i)=6-2=4>0,
\]
and
\[
\Delta_3=\det(H)=8>0.
\]
By Sylvester's criterion for Hermitian matrices, $H$ is positive definite. Hence, by the Hermitian positive definite matrix theorem,
$u^\dagger Hv$ defines an inner product.
\[
\boxed{\,\ip{\cdot}{\cdot}\text{ is an inner product on }\C^{3,1}.\,}
\]

\item Let $e_1,e_2,e_3$ be the standard basis. Because the Gram matrix is $H$, we can read off
\[
\ip{e_i}{e_j}=H_{ij}.
\]
Thus
\[
\ip{e_1}{e_1}=2,
\qquad
\ip{e_1}{e_2}=1+i,
\qquad
\ip{e_1}{e_3}=0.
\]
Using the projection formula
\[
\proj_{u}(v)=\frac{\ip{u}{v}}{\ip{u}{u}}u,
\]
we get
\[
u_1=e_1,
\qquad
u_2=e_2-\frac{\ip{u_1}{e_2}}{\ip{u_1}{u_1}}u_1
=e_2-\frac{1+i}{2}e_1.
\]
Next, since
\[
\ip{u_1}{e_3}=0,
\qquad
\ip{u_2}{u_2}=2,
\qquad
\ip{u_2}{e_3}=2,
\]
we obtain
\[
u_3=e_3-\frac{\ip{u_1}{e_3}}{\ip{u_1}{u_1}}u_1
       -\frac{\ip{u_2}{e_3}}{\ip{u_2}{u_2}}u_2
=e_3-u_2
=\frac{1+i}{2}e_1-e_2+e_3.
\]
Therefore
\[
\boxed{\,
(u_1,u_2,u_3)=\left(e_1,\ e_2-\frac{1+i}{2}e_1,\ \frac{1+i}{2}e_1-e_2+e_3\right).
\,}
\]

\item The same Gram-matrix calculations give
\[
\ip{u_1}{u_1}=\ip{u_2}{u_2}=\ip{u_3}{u_3}=2.
\]
Hence
\[
\boxed{\,
\beta=\left(\frac{u_1}{\sqrt2},\frac{u_2}{\sqrt2},\frac{u_3}{\sqrt2}\right)
\,}
\]
is orthonormal, namely
\[
\boxed{
\beta=
\left(
\frac{1}{\sqrt2}\!\begin{bmatrix}1\\0\\0\end{bmatrix},\ 
\frac{1}{\sqrt2}\!\begin{bmatrix}-\frac{1+i}{2}\\1\\0\end{bmatrix},\ 
\frac{1}{\sqrt2}\!\begin{bmatrix}\frac{1+i}{2}\\-1\\1\end{bmatrix}
\right).
}
\]

\item For $v=(1,0,i)^{\mathsf T}$, first compute
\[
Hv=\begin{bmatrix}2\\1+i\\4i\end{bmatrix}.
\]
Since $\beta$ is orthonormal, the Fourier coefficients are
\[
c_k=\ip{\beta_k}{v}=\beta_k^\dagger Hv.
\]
Substituting the three vectors gives
\[
c_1=\sqrt2,
\qquad
c_2=\sqrt2\,i,
\qquad
c_3=\sqrt2\,i.
\]
Therefore
\[
\boxed{\,(c_1,c_2,c_3)=\left(\sqrt2,\sqrt2\,i,\sqrt2\,i\right).\,}
\]
\end{enumerate}

This method is efficient because it avoids recomputing each inner product from scratch: the entries of $H$ already encode all pairwise inner products of the standard basis vectors.



\subsubsection{Method 3: Complete-square / Cholesky factorisation}
Let $u=(x,y,z)^{\mathsf T}$. A direct expansion gives
\[
u^{\dagger}Hu
=2|x|^2+(1+i)\overline{x}y+(1-i)\overline{y}x+3|y|^2+2\overline{y}z+2\overline{z}y+4|z|^2.
\]
This can be written as a sum of squares:
\[
\boxed{\,
u^{\dagger}Hu
=2\left|x+\frac{1+i}{2}y\right|^2+2|y+z|^2+2|z|^2.
\,}
\]
The right-hand side is nonnegative. It is zero only if
\[
z=0,
\qquad
y+z=0,
\qquad
x+\frac{1+i}{2}y=0.
\]
These equations imply $z=0$, then $y=0$, then $x=0$. Hence
\[
u^{\dagger}Hu>0\qquad(u\ne0).
\]
Since $H^{\dagger}=H$, the matrix $H$ is Hermitian positive definite. Therefore
\[
\ip{u}{v}=u^{\dagger}Hv
\]
defines an inner product.

Equivalently, define
\[
R\df
\begin{bmatrix}
1&\frac{1+i}{2}&0\\
0&1&1\\
0&0&1
\end{bmatrix}.
\]
Then
\[
H=2R^{\dagger}R.
\]
Thus
\[
\ip{u}{v}=2(Ru)^{\dagger}(Rv),
\]
so the given inner product is just a positive scalar multiple of the standard inner product after the invertible coordinate transformation $u\mapsto Ru$.

For Gram--Schmidt, use the already computed orthogonal basis
\[
u_1=e_1,
\qquad
u_2=e_2-\frac{1+i}{2}e_1,
\qquad
u_3=e_3-e_2+\frac{1+i}{2}e_1.
\]
The complete-square factorisation also shows efficiently that
\[
\|u_1\|^2=\|u_2\|^2=\|u_3\|^2=2.
\]
Hence the orthonormal basis is
\[
\boxed{\,
\beta=\left(\frac{u_1}{\sqrt2},\frac{u_2}{\sqrt2},\frac{u_3}{\sqrt2}\right).
\,}
\]
For $v=(1,0,i)^{\mathsf T}$, the Fourier coefficients are again
\[
\boxed{\,(c_1,c_2,c_3)=(\sqrt2,\sqrt2 i,\sqrt2 i).\,}
\]

\subsubsection{Method 4: One-line Gram--Schmidt from the Gram matrix}
Because the standard-basis Gram matrix is $H$, we have
\[
\ip{e_i}{e_j}=H_{ij}.
\]
Thus every needed inner product can be read directly from the entries of $H$.

Seek
\[
u_1=e_1,
\qquad
u_2=e_2-ae_1.
\]
The condition $u_2\perp u_1$ gives
\[
0=\ip{u_1}{u_2}=\ip{e_1}{e_2-ae_1}
=H_{12}-aH_{11}.
\]
Therefore
\[
a=\frac{H_{12}}{H_{11}}=\frac{1+i}{2},
\]
so
\[
u_2=e_2-\frac{1+i}{2}e_1.
\]
Now seek
\[
u_3=e_3-bu_1-cu_2.
\]
The projection coefficients are
\[
b=\frac{\ip{u_1}{e_3}}{\ip{u_1}{u_1}}=0,
\qquad
c=\frac{\ip{u_2}{e_3}}{\ip{u_2}{u_2}}=\frac{2}{2}=1.
\]
Hence
\[
u_3=e_3-u_2=e_3-e_2+\frac{1+i}{2}e_1.
\]
This produces the same answer as Method 1 while minimizing computation:
\[
\boxed{\,
(u_1,u_2,u_3)=\left(e_1,\ e_2-\frac{1+i}{2}e_1,\ e_3-e_2+\frac{1+i}{2}e_1\right).
\,}
\]


%====================================================
% QUESTION 5
%====================================================
\newpage
\section{Problem 5 (Double orthogonal complements)}
\subsection{Problem}
Let (i) $(V,\ip{\cdot}{\cdot})$ be an inner product space, and (ii) $W$ be a vector subspace of $V$.
\begin{enumerate}[label=(\alph*)]
\item Show that $W\subseteq W^{\perp\perp}$.
\item Show that if $V$ is finite-dimensional, then $W=W^{\perp\perp}$.
\item Show that if $V$ is infinite-dimensional, then it is not necessarily true that $W=W^{\perp\perp}$.
\end{enumerate}

\subsection{Solution}

\subsubsection{Method 1: Direct definition + dimension counting (finite-dimensional case)}
\begin{enumerate}[label=(\alph*)]
\item Let $w\in W$. To show $w\in W^{\perp\perp}$ we must show $\ip{w}{u}=0$ for all $u\in W^\perp$.
But $u\in W^\perp$ means precisely that $\ip{u}{w'}=0$ for all $w'\in W$; in particular $\ip{u}{w}=0$.
Using conjugate symmetry (or symmetry in the real case), $\ip{w}{u}=0$. Hence $w\in W^{\perp\perp}$.
So $W\subseteq W^{\perp\perp}$.

\[
\boxed{\,W\subseteq W^{\perp\perp}\,}
\]

\item Assume $\dim(V)<\infty$. A standard theorem gives
\[
\dim(W)+\dim(W^\perp)=\dim(V).
\]
Apply this to $W^\perp$:
\[
\dim(W^\perp)+\dim(W^{\perp\perp})=\dim(V).
\]
Subtracting yields $\dim(W^{\perp\perp})=\dim(W)$. Together with (a) and finite-dimensionality,
\[
W\subseteq W^{\perp\perp}\ \text{ and }\ \dim(W)=\dim(W^{\perp\perp})
\quad\Longrightarrow\quad
W=W^{\perp\perp}.
\]
\[
\boxed{\,\dim(V)<\infty\ \Rightarrow\ W=W^{\perp\perp}\,}
\]

\item In infinite dimensions, equality can fail. Let $V=\ell^2$ be the space of square-summable real sequences
(with the standard inner product $\ip{x}{y}=\sum_{k=1}^\infty x_k y_k$), and let
\[
W=c_{00}\df \{x\in\ell^2:\ x\text{ has finite support}\}.
\]
Then $W$ is a subspace of $V$ and is infinite-dimensional (but not all of $\ell^2$).

We claim $W^\perp=\{0\}$. Indeed, if $z\in W^\perp$, then for each standard basis vector $e_k\in c_{00}$ we have
\[
0=\ip{z}{e_k}=z_k,
\]
so all coordinates of $z$ are $0$, hence $z=0$.
Therefore $W^\perp=\{0\}$ and hence $W^{\perp\perp}=V$.

But $W\neq V$ (e.g.\ $(1,1/2,1/3,\dots)\in \ell^2$ but not finitely supported).
Thus $W\neq W^{\perp\perp}$.

\[
\boxed{\,\text{In infinite dimensions, }W\neq W^{\perp\perp}\text{ can occur (e.g.\ }W=c_{00}\subset \ell^2).\,}
\]
\end{enumerate}

\subsubsection{Method 2: Orthogonal decomposition (finite-dimensional case) + ``closure'' intuition}
\begin{enumerate}[label=(\alph*)]
\item Same as Method 1.

\item If $\dim(V)<\infty$, choose an orthonormal basis of $W$ and extend it to an orthonormal basis of $V$ by Gram--Schmidt:
\[
(\text{ONB of }W)\ \cup\ (\text{ONB of }W^\perp)=\text{ONB of }V.
\]
Then every $v\in V$ has a unique decomposition $v=w+u$ with $w\in W$, $u\in W^\perp$.
Now if $v\in W^{\perp\perp}$, then $\ip{v}{u}=0$ for all $u\in W^\perp$.
In particular, taking $u$ from the decomposition $v=w+u$ gives
\[
0=\ip{v}{u}=\ip{w+u}{u}=\ip{w}{u}+\ip{u}{u} = 0+\|u\|^2,
\]
so $u=0$ and hence $v=w\in W$. Thus $W^{\perp\perp}\subseteq W$, and with (a) we get equality.

\item In general (especially in infinite dimensions), $W^{\perp\perp}$ corresponds to the ``closed span'' of $W$ in the norm topology,
so if $W$ is not closed, strict inclusion may occur. The $\ell^2$ example above realizes this phenomenon concretely.
\end{enumerate}



\subsubsection{Method 3: Projection theorem proof in finite dimensions}
The key finite-dimensional theorem is the orthogonal decomposition theorem:
\[
V=W\oplus W^{\perp}.
\]
Thus every $v\in V$ has a unique decomposition
\[
v=w+u,
\qquad w\in W,
\qquad u\in W^{\perp}.
\]
We already know from Part (a) that
\[
W\subseteq W^{\perp\perp}.
\]
To prove the reverse inclusion in finite dimensions, let $v\in W^{\perp\perp}$. Write
\[
v=w+u,
\qquad w\in W,
\qquad u\in W^{\perp}.
\]
Since $v\in W^{\perp\perp}$, it is orthogonal to every vector in $W^{\perp}$, in particular to $u$. Therefore
\[
0=\ip{v}{u}=\ip{w+u}{u}=\ip{w}{u}+\ip{u}{u}.
\]
But $w\in W$ and $u\in W^{\perp}$, so $\ip{w}{u}=0$. Hence
\[
\|u\|^2=0,
\]
so $u=0$. Therefore
\[
v=w\in W.
\]
Thus
\[
W^{\perp\perp}\subseteq W.
\]
Together with $W\subseteq W^{\perp\perp}$, this gives
\[
\boxed{\,W=W^{\perp\perp}\quad\text{if }\dim V<\infty.\,}
\]
This proof is often more conceptual than dimension counting because it explicitly uses the perpendicular component of $v$.

\subsubsection{Method 4: Closure theorem in Hilbert spaces}
In a Hilbert space, the stronger theorem is
\[
\boxed{\,W^{\perp\perp}=\overline{W}.\,}
\]
Here $\overline{W}$ is the closure of $W$ under the norm induced by the inner product.

This theorem explains both the finite-dimensional and infinite-dimensional cases at once. In finite dimensions, every subspace is closed, so
\[
\overline{W}=W,
\]
and hence
\[
W^{\perp\perp}=W.
\]
In infinite dimensions, a subspace need not be closed. If $W$ is dense but not equal to $V$, then
\[
\overline{W}=V,
\]
so
\[
W^{\perp\perp}=V\ne W.
\]
For the example in Method 1,
\[
W=c_{00}\subset \ell^2.
\]
The subspace $c_{00}$ of finitely supported sequences is dense in $\ell^2$, so
\[
\overline{c_{00}}=\ell^2.
\]
Therefore
\[
c_{00}^{\perp\perp}=\ell^2\ne c_{00}.
\]
This is the advanced conceptual reason why Part (c) is possible.


%====================================================
% QUESTION 6
%====================================================
\newpage
\section{Problem 6 (Orthogonal complement and orthogonal decomposition in $\R^4$)}
\subsection{Problem}
Let
\[
V\df \{(w,x,y,z)\in\R^4\mid w-y+z=0\},
\]
which is a vector subspace of $\R^4$ (standard inner product $\cdot$).
\begin{enumerate}[label=(\alph*)]
\item Derive a basis for $V^\perp$.
\item Derive an orthonormal basis for $(V,\cdot)$.
\item Compute the unique $a\in V$ and $b\in V^\perp$ such that $(1,2,3,4)=a+b$.
\end{enumerate}

\subsection{Solution}

\subsubsection{Method 1: Solve linear constraints for $V^\perp$; then Gram--Schmidt inside $V$}
\begin{enumerate}[label=(\alph*)]
\item Parameterize $V$ by free variables $x,y,z$:
\[
w=y-z,\qquad (w,x,y,z)=(y-z,\ x,\ y,\ z)=x(0,1,0,0)+y(1,0,1,0)+z(-1,0,0,1).
\]
Let $u=(a,b,c,d)\in\R^4$. Then $u\in V^\perp$ iff $u\cdot v=0$ for all $v\in V$.
It suffices to impose orthogonality to the spanning vectors:
\[
u\cdot(0,1,0,0)=b=0,\qquad
u\cdot(1,0,1,0)=a+c=0,\qquad
u\cdot(-1,0,0,1)=-a+d=0.
\]
Thus $b=0$, $c=-a$, $d=a$, so
\[
V^\perp=\Span\{(1,0,-1,1)\}.
\]
\[
\boxed{\,V^\perp=\Span\{(1,0,-1,1)\}\,}.
\]

\item A convenient basis of $V$ from the parameterization is
\[
v_1=(0,1,0,0),\quad v_2=(1,0,1,0),\quad v_3=(-1,0,0,1).
\]
Apply Gram--Schmidt:

\underline{Step 1:} $e_1=v_1$ (already unit), so $e_1=(0,1,0,0)$.

\underline{Step 2:} $v_2\cdot e_1=0$, so $u_2=v_2$ and
\[
e_2=\frac{v_2}{\|v_2\|}=\frac{1}{\sqrt2}(1,0,1,0).
\]

\underline{Step 3:} Orthogonalize $v_3$ against $u_2=v_2$:
\[
v_3\cdot v_2=(-1)\cdot 1+0\cdot 0+0\cdot 1+1\cdot 0=-1,\qquad v_2\cdot v_2=2.
\]
So
\[
u_3=v_3-\proj_{v_2}(v_3)=v_3-\frac{-1}{2}v_2=v_3+\frac12 v_2=\left(-\frac12,0,\frac12,1\right).
\]
Equivalently, $2u_3=(-1,0,1,2)$, and $\|(-1,0,1,2)\|^2=6$. Hence
\[
e_3=\frac{1}{\sqrt6}(-1,0,1,2).
\]
Thus an orthonormal basis for $V$ is
\[
\boxed{\,
\left(
(0,1,0,0),\ \frac{1}{\sqrt2}(1,0,1,0),\ \frac{1}{\sqrt6}(-1,0,1,2)
\right).
\,}
\]

\item Let $p=(1,2,3,4)$. Since $V^\perp=\Span\{u\}$ with $u=(1,0,-1,1)$, we have
\[
b=\proj_{V^\perp}(p)=\frac{p\cdot u}{u\cdot u}u.
\]
Compute
\[
p\cdot u=1\cdot 1+2\cdot 0+3\cdot(-1)+4\cdot 1=2,\qquad u\cdot u=1+0+1+1=3.
\]
So
\[
b=\frac{2}{3}(1,0,-1,1)=\left(\frac23,0,-\frac23,\frac23\right).
\]
Then
\[
a=p-b=\left(1-\frac23,\ 2,\ 3+\frac23,\ 4-\frac23\right)=\left(\frac13,\ 2,\ \frac{11}{3},\ \frac{10}{3}\right).
\]
Check $a\in V$:
\[
w-y+z=\frac13-\frac{11}{3}+\frac{10}{3}=0.
\]
Therefore
\[
\boxed{\,
a=\left(\frac13,\ 2,\ \frac{11}{3},\ \frac{10}{3}\right)\in V,
\quad
b=\left(\frac23,0,-\frac23,\frac23\right)\in V^\perp,
\quad
(1,2,3,4)=a+b.
\,}
\]
\end{enumerate}

\subsubsection{Method 2: Matrix viewpoint (constraint normal vector)}
The subspace $V=\{x\in\R^4: n\cdot x=0\}$ where $n=(1,0,-1,1)$ (since $w-y+z=0$).
Hence $V^\perp=\Span\{n\}$ immediately. The orthogonal decomposition of $p$ is
\[
p=\left(p-\frac{p\cdot n}{n\cdot n}n\right)+\frac{p\cdot n}{n\cdot n}n,
\]
matching the $a,b$ above.



\subsubsection{Method 3: Kernel--row-space theorem}
Write the defining equation of $V$ as a matrix equation:
\[
w-y+z=0
\quad\Longleftrightarrow\quad
\begin{bmatrix}1&0&-1&1\end{bmatrix}
\begin{bmatrix}w\\x\\y\\z\end{bmatrix}=0.
\]
Hence
\[
V=\Null(A),
\qquad
A\df\begin{bmatrix}1&0&-1&1\end{bmatrix}.
\]
A standard theorem says
\[
\Null(A)^{\perp}=\Row(A).
\]
Therefore
\[
V^{\perp}=\Row(A)=\Span\{(1,0,-1,1)\}.
\]
So
\[
\boxed{\,V^{\perp}=\Span\{(1,0,-1,1)\}.\,}
\]
The orthogonal decomposition of $p=(1,2,3,4)$ then follows by projecting onto the row vector
\[
n=(1,0,-1,1).
\]
Thus
\[
b=\frac{p\cdot n}{n\cdot n}n
=\frac{2}{3}(1,0,-1,1)
=\left(\frac23,0,-\frac23,\frac23\right),
\]
and
\[
a=p-b=\left(\frac13,2,\frac{11}{3},\frac{10}{3}\right).
\]

\subsubsection{Method 4: Lagrange multiplier / closest-point method}
The vector $a\in V$ in the decomposition
\[
p=a+b,
\qquad a\in V,
\qquad b\in V^{\perp},
\]
is the closest point in the hyperplane $V$ to $p=(1,2,3,4)$.
Thus we minimize
\[
\|p-a\|^2
\]
subject to the constraint
\[
n\cdot a=0,
\qquad n=(1,0,-1,1).
\]
The closest-point condition says that $p-a$ must be normal to the hyperplane, so
\[
p-a=\lambda n.
\]
Hence
\[
a=p-\lambda n.
\]
Impose $a\in V$:
\[
0=n\cdot a=n\cdot(p-\lambda n)=n\cdot p-\lambda n\cdot n.
\]
Therefore
\[
\lambda=\frac{n\cdot p}{n\cdot n}=\frac{2}{3}.
\]
So
\[
b=p-a=\lambda n=\frac23(1,0,-1,1),
\]
and
\[
a=p-b=\left(\frac13,2,\frac{11}{3},\frac{10}{3}\right).
\]
This method is especially useful for projection questions involving one linear constraint.

\subsubsection{Method 5: Constructing the orthonormal basis by inspection}
Since
\[
V=n^{\perp},
\qquad n=(1,0,-1,1),
\]
we need three mutually orthogonal unit vectors perpendicular to $n$.
One immediate choice is
\[
e_1=(0,1,0,0),
\]
because its dot product with $n$ is $0$.
Now look only at the $(w,y,z)$ coordinates satisfying
\[
w-y+z=0.
\]
Two convenient perpendicular choices are
\[
v_2=(1,0,1,0),
\qquad
v_3=(-1,0,1,2).
\]
Indeed,
\[
v_2\cdot n=1-1=0,
\qquad
v_3\cdot n=-1-1+2=0,
\]
and
\[
v_2\cdot v_3=-1+1=0.
\]
Also both are perpendicular to $e_1$. Their norms are
\[
\|e_1\|=1,
\qquad
\|v_2\|=\sqrt2,
\qquad
\|v_3\|=\sqrt6.
\]
Thus an orthonormal basis is
\[
\boxed{\,
\left(
(0,1,0,0),\ \frac{1}{\sqrt2}(1,0,1,0),\ \frac{1}{\sqrt6}(-1,0,1,2)
\right).
\,}
\]
This gives the same basis as Method 1 but avoids running the full Gram--Schmidt algorithm.


%====================================================
% QUESTION 7
%====================================================
\newpage
\section{Problem 7 (Diagonal matrices and Frobenius orthogonal complement)}
\subsection{Problem}
Let $n\in\mathbb{N}$, and consider the real inner product space $(\R^{n,n},\ip{\cdot}{\cdot}_F)$, where
\[
\ip{A}{B}_F\df \tr(A^{\mathsf T}B)=\sum_{i=1}^n\sum_{j=1}^n a_{ij}b_{ij}.
\]
Let
\[
D\df \{A\in\R^{n,n}\mid A\ \text{is diagonal}\}.
\]
Derive a basis for $D^\perp$.

\subsection{Solution}

\subsubsection{Method 1: Characterize $D^\perp$ by testing against diagonal matrices}
Let $X=(x_{ij})\in\R^{n,n}$. Then $X\in D^\perp$ iff $\ip{X}{\Delta}_F=0$ for all diagonal $\Delta=\diag(d_1,\dots,d_n)$.
But
\[
\ip{X}{\Delta}_F=\sum_{i=1}^n\sum_{j=1}^n x_{ij}\Delta_{ij}=\sum_{i=1}^n x_{ii}d_i.
\]
For this to be $0$ for all choices of $(d_1,\dots,d_n)\in\R^n$, we must have $x_{ii}=0$ for every $i$.
Thus
\[
D^\perp=\{X\in\R^{n,n}: x_{11}=x_{22}=\cdots=x_{nn}=0\},
\]
i.e.\ the subspace of matrices with zero diagonal.

Let $E_{ij}$ be the standard matrix units. Then $\{E_{ij}: i\neq j\}$ spans precisely the zero-diagonal matrices and is linearly independent.
Therefore
\[
\boxed{\,D^\perp=\Span\{E_{ij}: i\neq j\}\quad\text{and}\quad \{E_{ij}: i\neq j\}\ \text{is a basis of }D^\perp.\,}
\]

\subsubsection{Method 2: Orthogonal direct sum decomposition}
Every matrix $A\in\R^{n,n}$ decomposes uniquely as
\[
A=\underbrace{\diag(a_{11},\dots,a_{nn})}_{\in D}\ +\ \underbrace{\bigl(A-\diag(a_{11},\dots,a_{nn})\bigr)}_{\text{zero diagonal}}.
\]
Moreover, diagonal matrices are Frobenius-orthogonal to zero-diagonal matrices, since
\[
\ip{\diag(d_1,\dots,d_n)}{X}_F=\sum_{i=1}^n d_i x_{ii}=0
\]
whenever $x_{ii}=0$. Hence $\R^{n,n}=D\oplus D^\perp$ and $D^\perp$ is exactly the zero-diagonal subspace,
with basis $\{E_{ij}:i\neq j\}$.


\subsubsection{Method 3: Matrix-unit orthonormal basis}
Let $E_{ij}$ denote the standard matrix unit whose $(i,j)$-entry is $1$ and whose other entries are $0$.
Under the Frobenius inner product,
\[
\ip{E_{ij}}{E_{kl}}_F
=\tr(E_{ij}^{\mathsf T}E_{kl})
=\delta_{ik}\delta_{jl}.
\]
Therefore
\[
\{E_{ij}:1\le i,j\le n\}
\]
is an orthonormal basis of $\R^{n,n}$.
The diagonal subspace is
\[
D=\Span\{E_{11},E_{22},\dots,E_{nn}\}.
\]
Since the full matrix-unit basis is orthonormal, the orthogonal complement of $D$ is spanned by all matrix units not appearing in the basis for $D$:
\[
\boxed{\,D^{\perp}=\Span\{E_{ij}:i\ne j\}.\,}
\]
This is the fastest method because it identifies $D$ as a coordinate subspace of the Frobenius matrix space.

\subsubsection{Method 4: Orthogonal projection onto the diagonal subspace}
Define the diagonal projection
\[
\Pi_D(A)\df \diag(a_{11},a_{22},\dots,a_{nn}).
\]
Then every matrix $A\in\R^{n,n}$ decomposes uniquely as
\[
A=\Pi_D(A)+\bigl(A-\Pi_D(A)\bigr).
\]
The first term is diagonal, and the second term has zero diagonal. Moreover, if $\Delta=\diag(d_1,\dots,d_n)$ and $X$ has zero diagonal, then
\[
\ip{\Delta}{X}_F=\sum_{i=1}^{n}d_i x_{ii}=0.
\]
Hence the decomposition is an orthogonal direct sum:
\[
\R^{n,n}=D\oplus D^{\perp}.
\]
Therefore
\[
D^{\perp}=\ker(\Pi_D)=\{X\in\R^{n,n}:x_{11}=x_{22}=\cdots=x_{nn}=0\}.
\]
A basis is again
\[
\boxed{\,\{E_{ij}:i\ne j\}.\,}
\]


\end{document}
